\section{Experimental protocol}
\label{sec:setup}
\subsection{Development diagnostics and frozen final tasks}
The ChangeOS camera-ladder diagnostic uses 43 development ROIs containing 7,238 annotated buildings. It compares cruise, intermediate, and floor predictions on the same ground-truth building identities. The full-ground-truth endpoint counts an unmatched building as incorrect; matched-only classification and probability metrics use only matched detections. These are development and mechanism measurements, not final policy scores.

The confirmatory online question set is separate from the 160-question development set used to choose policy rules. It contains 160 questions from 44 previously unconsumed ROIs: 48 in Hurricane Michael, 51 in the Moore tornado, and 61 in Nepal flooding. Questions cover presence, damage, count, and spatial relations. The author manually reviewed all 160 and approved the frozen set; 88 previously flagged ambiguity notices remain recorded rather than erased. The present record documents one author's review, not two independent annotators or an inter-rater agreement estimate. The final question set, review report, binary-label manifest, and source/prompt/weight hashes were checked before the first online episode. The exact identities and checksums are listed in Appendix~\ref{app:provenance}.

\subsection{Online execution and endpoints}
Seven policies run on each of the 160 final questions for each of three generation repeats, producing 3,360 planned episodes. Four shards partition each repeat without changing the question set. The action base seed is 42; generation calls derive their seed from base 42000 and $(\mathrm{qid},\mathrm{repeat},\mathrm{step},\mathrm{call\ role})$. Qwen sampling settings, source fingerprint, prompt hash, and ChangeOS weights are frozen in the protocol. All 12 shard audits passed with zero execution errors. Repeats provide three generation conditions on the same 160 questions; they are not new independent questions, ROIs, or events.

The prespecified primary endpoint is terminal task accuracy over all questions, averaged over repeats. Abstentions count as incorrect for this endpoint. Secondary descriptions include accuracy by event and task type, abstention, within-policy answer correction and harm, executed reobservations, state-derived horizontal/vertical motion, and episode wall time. A correction or harm observed inside one policy's trajectory is not the counterfactual effect of its action relative to another policy. Later VLM call counts can also differ across policies, so paired accuracy contrasts estimate whole-policy outcomes rather than a pure image-change effect.

Four T1\_TASK contrasts against A0\_HOLD, A1\_RANDOM, A2\_ALWAYS, and A3U\_RAW\_ENTROPY were declared before execution. Pairing uses the same $(\mathrm{repeat},\mathrm{qid})$ units. For each contrast, a 5,000-draw ROI-cluster bootstrap resamples all questions and three repeats within each sampled ROI, with seed 42. The four primary comparisons use Bonferroni-adjusted 98.75\% percentile intervals; 95\% intervals are descriptive. Only three events are represented, so these intervals characterize variation among the observed ROI clusters, not transfer to unseen disasters. Event and question-type breakdowns are descriptive rather than additional confirmatory tests.

Executed reobservation count is the principal observed action-cost measure. Episode wall time is measured software runtime for this execution, not UAV flight time. Trajectory-state displacements are geometric estimates, and configured-speed $d/v$ values are estimates rather than measured airborne duration. No randomized common-cap sweep or component-level render/ChangeOS/controller/Qwen timing is available. Accordingly, accuracy and cost are reported as realized policy operating points.
